Una flauta travessera és un tub metàl·lic obert pels dos extrems que té una longitud de 67,0~cm. El seu so abraça un interval extens de freqüències i és un instrument molt comú en orquestres.

\begin{apartat}{1,25}
Per a tocar la nota més greu, el flautista ha de tapar amb els dits tots els forats laterals del tub. Calculeu les freqüències del primer i el tercer harmònics d'una flauta travessera en aquest cas. Dibuixeu aquests dos harmònics i calculeu per a cada un la posició dels nodes i els ventres respecte d'un extrem de la flauta.
\end{apartat}

\begin{apartat}{1,25}
Tapem un dels extrems de la flauta de manera que aquesta es comporta com un tub amb un extrem obert i un de tancat. Dibuixeu l'ona estacionària corresponent al primer i el segon modes de vibració possibles i indiqueu els seus nodes i ventres. Calculeu la longitud d'ona i la freqüència d'aquests dos modes de vibració.
\end{apartat}

\dades{Velocitat del so en l'aire $= 343\ \mathrm{m\,s^{-1}}$.}
